\noindent\textit{Where the pixels change.} Write $\rho=\sum_{p\in M_{\rm edit}}\delta_p/\sum_p\delta_p$, where $\delta_p$ is mean absolute RGB change and $M_{\rm edit}$ is the edit mask. The background receives fraction $1-\rho$ of total change. On the 42 seed-0 SAM images, the constants lower whole-image MAE from 0.0600 to 0.0344 while the edit share stays close (uniform 0.380, constants 0.371). Spatial lowers MAE to 0.0283 and raises the share to 0.480. The constants chiefly reduce the amount of change. Spatial weighting also concentrates it inside the edit mask. In the restoring study at $U$, mean total change is almost identical across allocations while spatial raises the edit share by 0.015 [0.008, 0.023].

Pixel selectivity $S_{\rm MAE}$ is the mean of per-image edit/background MAE ratios. Across uniform settings whole-image change grows nearly eightfold but selectivity stays within 1.88--2.11; the constants give 2.08. Spatial selectivity rises from 2.07 at scale 0.25 to 4.00 at the default and 4.88 at the extended scale. At scale 1.25 the two arms have similar cohort-mean directional CLIP (0.1239, 0.1235) and different background LPIPS (0.1033, 0.3268). Equation~\eqref{eq:mrc} addresses within-image response differences.

With scales up to one, uniform has lower matched background LPIPS for the jumping dog, the square crate made round and the rice field. With all five settings, the exceptions are the jumping dog, the reflection, the jar lid and the rice field. The jar favors spatial at default on both background LPIPS (0.030 versus 0.045) and directional CLIP (0.048 versus 0.014). Its response ranges first overlap when scale 1.25 is included, giving $U-S=-0.0514$ in the five-setting comparison. Across the four supported pose edits in that scan, the mean changes sign with the scorer ($+0.0051$ FM, $-0.0013$ masked input). A changed outline may extend outside the source mask.
